%!TEX root = main.tex

\section{Introduction}

With the rise of the Internet of Things (IoT), both the number and the diversity of connected devices is staggering.
%
We have entered in a new era where the objects of our daily life, from mobile to wearable devices, can be interconnected.
%
Leaders of the industry, such as Cisco, predict that 50 billion devices will be connected together by 2020.
%
The goal of IoT is not really to enable users to interact with everything.
%
Instead, it is more about cross-linking objects to promote machine-to-machine interactions.
%
These interactions enable the capture and/or distribution of data that can be turned into information and delivered to end-users.




IoT has led to a strong revival of interest in mobile ad-hoc networks (MANET).
%
A MANET is a mobile network without a fixed infrastructure; meaning that it is self-configuring network where nodes connect directly to each other.
%
In addition, the topology of the network changes frequently due to the mobility of the participants, nodes are connected using a wireless medium, there is no central entity in charge of routing, and most participants operate on batteries with limited processing power and bandwidth.
%
Applications based on MANETs range from sensor networks to end-user mobile applications, such as traffic jam prevention system, information dissemination in crowds, data gathering in hostile environments, and peer-to-peer mobile games.



Many network functions in MANETs, such as routing for instance, rely on the fundamental \textit{broadcast operation}, which consists in delivering a message from a source node to all other nodes in a network.
%
Since the transmission range of a node is usually smaller than the size of the network, the solution is to have some nodes retransmiting the message to reach all the network.
%
The broadcasting problem in MANET have received considerable attention.
%
In particular, the solutions attempt to:

\begin{itemize}
	\item \textit{Maximize coverage.} As many nodes in the network as possible receive a broadcast message $m$.
	\item \textit{Minimize energy consumption and bandwidth usage.} Ultimately, the efficiency of an algorithm is measured by the amount of energy, processing power and network bandwidth used to disseminate messages.
	\item \textit{Support mobility patterns.} The velocity, pauses and number of nodes in a location impact both coverage and efficiency. For that reason, it is useful to consider particular deployment scenarios. 
\end{itemize}



Many algorithms have been proposed to solve the broadcasting problem in MANET~\cite{}.
%
They seek to control the flooding by selecting which nodes should transmit a message during dissemination and when these nodes should transmit.
%
Algorithms' performance depend on their ability to use knowledge about the network topology, node location, and network status.
%
However, their performance is also largely influenced by external factors such as the density of the network, and the mobility rate of its nodes.
%
Although the performance of previous solutions have been studied and extensive work have been done in classifying algorithms based on their characteristic, to the best of our knowledge, there is a lack of comprehensive comparison of the existent solutions.  
%
We propose to explore the relative performance and costs of representative MANET broadcast algorithms in a variety of deployment conditions.



Our study considers five protocols, in addition to the simple flooding algorithm, that are representative of their class:

\er{This is a proposal for an introduction structure, to be discussed}

\begin{itemize}
%	\item Definition and importance of MANETs
%	\item Characteristics of MANETs
%	\begin{itemize}
%		\item no access to wired network connection
%		\item use a wireless communication medium
%		\item operate on battery
%		\item limited processing power
%		\item mobile in space
%	\end{itemize}
%	\item Operating MANETs rely on a set of fundamental operations, among which broadcast is of great importance \er{ideally we can link to work that relies on broadcast to form routing paths}
%	\item Challenges in MANET networks: limited energy, mobility, wireless communication and collisions, limited knowledge of the topology, etc.
%	\item Broadcast received considerable attention in order to maximize coverage, minimize cost, support mobility patterns, and minimize the consumption of energy
%	\item Protocols in the literature seek to control the flooding of the broadcast message by selecting which node should transmit wirelessly, and when:
%	\begin{itemize}
%		\item Some protocols control the diffusion protocol in space, determining either in advance (topology-based) or during the dissemination, which node is eligible to forward the message when it receives it
%		\item Other protocols control the diffusion protocol in time, determining when a received message should be forwarded and under which conditions, avoiding unnecessary collisions and retransmissions
%	\end{itemize}
%	\item Protocol efficiency depends on their ability to use knowledge about the network topology, node location, and 
%	\item But the efficiency of protocols is also largely influenced by external factors, and primarily by the density of the network, and the mobility rate of its nodes. \er{not sure we should mention other aspects, such as the loosiness of channels, the presence of obstacles, etc.?}
%	\item The relative efficiency of each of the broadcast protocols against the simple and inefficient flooding algorithm, where no control is performed on the flooding operation, has been studied but the relative performance of advanced protocols between them is largely unstudied
%	\item We propose to explore the relative performance and costs of representative MANET broadcast algorithms in a variety of deployment conditions.
	\item Our study considers three \er{more?} protocols, in addition to the baseline flooding algorithm, that are representative of their class:
	\begin{itemize}
		\item ABBA controls the dissemination process in time, by setting up timers prior to dissemination that are inversely proportional to the probability that the message will be received from its neighbors from other sources
		\item CDS builds and maintain an overlay independently from dissemination, allowing to control the latter in space by selecting which node should or should not propagate incoming messages
		\item MPR is another well-known algorithms that maintains an overlay structure to control the broadcasting process in space. This overlay is built using a heuristic that aims at minimizing the number of relay nodes required to reach full coverage.
	\end{itemize}
	\item Our study allows answering the question of what is the most appropriate protocol for a certain deployment scenario, and concludes that there is no one-size-fits-all in MANET broadcasting \er{provide summary of main findings}
\end{itemize}